
\vspace{-0.25cm}
% Dans la partie « Dynamique dans un référentiel non galiléen », l'étude du champ de pesanteur est conduite en supposant le référentiel géocentrique galiléen. De nombreuses applications permettent d'illustrer cette partie : le pendule de Foucault, la déviation vers l'est, les vents géostrophiques, les courants marins ; l'étude statique des marées constitue également une ouverture pertinente.
{\footnotesize
\begin{tabularx}{\textwidth}[t]{|X|X|}
  \hline 
  \textbf{Notions et contenus} & \textbf{Capacités exigibles} \\
  \hline
  \multicolumn{2}{|l|}{\text { 4.3.4. Bilans macroscopiques }} \\
  \hline 
  Bilans de masse.  & Établir un bilan de masse en raisonnant sur un système ouvert et fixe ou sur un système fermé et mobile. \\ 
  \hline
  Bilans de quantité de mouvement ou d'énergie cinétique pour un écoulement stationnaire unidimensionnel à une entrée et une sortie. & Bilans de quantité de mouvement ou d'énergie cinétique pour un écoulement stationnaire unidimensionnel à une entrée et une sortie. \\
  \hline
\end{tabularx}
}